\begin{proof}
We argue by induction on \(n\).

\textbf{Base case \(n=1\).}  
If \(\lambda_{k}=1\) for some index \(k\) and \(\lambda_{i}=0\) for \(i\neq k\), then
\[
\sum_{i=1}^{n}\lambda_{i}x_{i}=x_{k},\qquad 
\sum_{i=1}^{n}\lambda_{i}f(x_{i})=f(x_{k}),
\]
so Jensen’s inequality
\[
f\!\Big(\sum_{i=1}^{n}\lambda_{i}x_{i}\Big)\le\sum_{i=1}^{n}\lambda_{i}f(x_{i}) \tag{J}
\]
holds with equality. This also covers the situation where some \(\lambda_{i}=0\).

\textbf{Binary convexity.}  
For two points \(x,y\in I\) and a weight \(t\in[0,1]\) the defining property of convexity,
\[
f\big(t x+(1-t) y\big)\le t f(x)+(1-t)f(y) \tag{C}
\]
is the elementary case from which the general inequality will be built.

\textbf{Induction hypothesis.}  
Assume that for a fixed integer \(k\ge 2\) the statement (J) holds for any collection of
\(k\) points and any non‑negative weights summing to \(1\); i.e.
\[
f\!\Big(\sum_{i=1}^{k}\lambda_{i}x_{i}\Big)\le\sum_{i=1}^{k}\lambda_{i}f(x_{i}) \tag{IH\(_k\)}
\]
whenever \(\lambda_{i}\ge0\) and \(\sum_{i=1}^{k}\lambda_{i}=1\).

\textbf{Inductive step \(k\to k+1\).}  
Let \(\lambda_{1},\dots ,\lambda_{k+1}\ge0\) with \(\sum_{i=1}^{k+1}\lambda_{i}=1\).  
Set
\[
\mu:=\lambda_{k+1},\qquad t:=1-\mu=\sum_{i=1}^{k}\lambda_{i}\;( \ge 0).
\]
If \(t=0\) then \(\mu=1\) and we are in the base case, so (J) holds with equality.
Hence we may assume \(t>0\).

Define the normalized convex combination of the first \(k\) points:
\[
y:=\frac{1}{t}\sum_{i=1}^{k}\lambda_{i}x_{i}.
\]
Because the coefficients \(\lambda_{i}/t\) are non‑negative and sum to \(1\), the induction
hypothesis applied to the \(k\) points \(x_{1},\dots ,x_{k}\) yields
\[
f(y)\le\frac{1}{t}\sum_{i=1}^{k}\lambda_{i}f(x_{i}). \tag{2}
\]

Now apply the binary convexity inequality \((\text{C})\) to the pair \((y,x_{k+1})\) with
weight \(t\):
\[
f\big(t y+\mu x_{k+1}\big)\le t f(y)+\mu f(x_{k+1}). \tag{3}
\]

\textit{Algebraic identity.}  
Since \(y=\frac{1}{t}\sum_{i=1}^{k}\lambda_{i}x_{i}\),
\[
t y+\mu x_{k+1}= \sum_{i=1}^{k}\lambda_{i}x_{i}+\lambda_{k+1}x_{k+1}
= \sum_{i=1}^{k+1}\lambda_{i}x_{i}. \tag{4}
\]

Substituting \((2)\) into the right–hand side of \((3)\) gives
\[
\begin{aligned}
f\!\Big(\sum_{i=1}^{k+1}\lambda_{i}x_{i}\Big)
&\le t\Big(\frac{1}{t}\sum_{i=1}^{k}\lambda_{i}f(x_{i})\Big)+\mu f(x_{k+1}) \\
&= \sum_{i=1}^{k}\lambda_{i}f(x_{i})+\lambda_{k+1}f(x_{k+1}) \\
&= \sum_{i=1}^{k+1}\lambda_{i}f(x_{i}),
\end{aligned}
\]
which is precisely (J) for \(k+1\) points. Thus the inductive step is proved.

\textbf{Completion of induction.}  
Since the inequality holds for the base case \(n=1\) and the inductive step is valid
for every \(k\ge 1\), by the principle of mathematical induction inequality (J) holds
for all integers \(n\ge 1\).

\textbf{Edge cases.}  
If some \(\lambda_{i}=0\), the corresponding term disappears from both sides; the
remaining weights still sum to \(1\), so the argument above applies unchanged.
If exactly one weight equals \(1\), the inequality reduces to equality, as shown in
the base case.

\textbf{Equality condition.}  
Equality in (J) can occur only when every inequality used in the proof becomes an
equality.

\begin{itemize}
\item In the binary convexity \((\text{C})\) equality holds iff either \(t\in\{0,1\}\) (i.e. one
weight is zero) or \(f\) is affine on the line segment joining the two points. For a
strictly convex \(f\) the latter forces the two points to coincide.
\item In the induction hypothesis \((\text{IH}_k)\) equality holds exactly when the same
condition holds for the reduced set of points.
\end{itemize}
Putting these together, equality in Jensen’s inequality holds \emph{iff} one (or more)
of the following is true:

\begin{enumerate}
\item \textbf{Weight concentration:} there exists \(j\) with \(\lambda_{j}=1\) and
\(\lambda_{i}=0\) for \(i\neq j\);
\item \textbf{All points equal:} \(x_{1}=x_{2}= \dots =x_{n}\) (then the convex
combination equals each \(x_{i}\));
\item \textbf{Affineness on the convex hull:} the function \(f\) is affine on the
convex hull of \(\{x_{1},\dots ,x_{n}\}\). In that case the inequality reduces to an
equality for any set of weights.
\end{enumerate}
For a strictly convex \(f\) the only possibilities are (1) and (2).

\end{proof}